\documentclass[aspectratio=169]{beamer}
\input{header}
\coursecode{JKSSB}

\title{Digital Platforms \& E-Governance}
\subtitle{Lesson 57 --- Named Platforms and G2x Recognition}
\author{JKSSB Computer Awareness}
\date{\today}

\begin{document}
\frame{\titlepage}

\begin{frame}{What e-Governance Means}
  \begin{itemize}
    \item \key{e-Governance} is the use of Information and Communication
      Technology (ICT) by government to deliver services.
    \item It makes services accessible, efficient, transparent, and reliable.
    \item It shifts service delivery from physical counters to online portals.
    \item e-Governance is about \key{governance}, not merely technology.
  \end{itemize}
  \begin{block}{The one idea}
    e-Governance uses ICT so that government services reach the common man in
    his locality, at affordable cost.
  \end{block}
\end{frame}

\begin{frame}{The G2x Family}
  \begin{center}
    \begin{tabular}{p{0.14\textwidth}p{0.34\textwidth}p{0.40\textwidth}}
      \toprule
      \textbf{Type} & \textbf{Full form} & \textbf{Who is served} \\
      \midrule
      G2C & Government to Citizen & The general public \\
      G2G & Government to Government & Government departments and levels \\
      G2B & Government to Business & Businesses and companies \\
      G2E & Government to Employee & Government employees \\
      \bottomrule
    \end{tabular}
  \end{center}
  \vspace{0.25cm}
  \muted{The letters say who talks to whom: C = citizen, G = government,
  B = business, E = employee.}
\end{frame}

\begin{frame}{Recognising G2x in Examples}
  \begin{itemize}
    \item \key{G2C}: UMANG, DigiLocker, applying for a certificate, paying a
      utility bill.
    \item \key{G2G}: e-Office, PFMS (working between government
      departments).
    \item \key{G2B}: GeM (Government e-Marketplace), GSTN.
    \item \key{G2E}: an employee service book or salary and leave portal.
  \end{itemize}
  \begin{alertblock}{Trap alert}
    GeM (government procurement) is G2B, not G2C; e-Office is G2G, not G2C.
  \end{alertblock}
\end{frame}

\begin{frame}{Digital India}
  \begin{itemize}
    \item \key{Digital India} is a flagship programme of the Government of
      India.
    \item Launched on \key{1 July 2015} by the Prime Minister.
    \item Vision: to transform India into a \key{digitally empowered society
      and knowledge economy}.
  \end{itemize}
  \begin{block}{Three vision areas}
    (i) Digital Infrastructure as a Core Utility to Every Citizen;
    (ii) Governance and Services on Demand;
    (iii) Digital Empowerment of Citizens.
  \end{block}
\end{frame}

\begin{frame}{Nine Pillars of Digital India}
  \begin{center}
    \begin{tabular}{p{0.06\textwidth}p{0.72\textwidth}}
      \toprule
      \textbf{\#} & \textbf{Pillar} \\
      \midrule
      1 & Broadband Highways \\
      2 & Universal Access to Mobile Connectivity \\
      3 & Public Internet Access Programme \\
      4 & e-Governance: Reforming Government through Technology \\
      5 & e-Kranti: Electronic Delivery of Services \\
      6 & Information for All \\
      7 & Electronics Manufacturing \\
      8 & IT for Jobs \\
      9 & Early Harvest Programmes \\
      \bottomrule
    \end{tabular}
  \end{center}
  \vspace{0.2cm}
  \muted{Pillar 5, e-Kranti, is the service-delivery pillar.}
\end{frame}

\begin{frame}{e-Kranti}
  \begin{itemize}
    \item \key{e-Kranti} means ``electronic delivery of services''.
    \item It is \key{Pillar 5} of the Digital India programme.
    \item It is also called \key{NeGP 2.0} --- an upgraded approach to the
      National e-Governance Plan.
    \item Vision: ``\key{Transforming e-Governance for Transforming
      Governance}''.
    \item Mission: deliver all government services electronically through
      integrated and interoperable systems.
  \end{itemize}
\end{frame}

\begin{frame}{NeGP --- National e-Governance Plan}
  \begin{itemize}
    \item \key{NeGP} = National e-Governance Plan.
    \item Approved in \key{May 2006}.
    \item Vision: make all government services accessible to the common man
      in his locality, through common service delivery outlets, at affordable
      cost.
    \item Built around \key{Mission Mode Projects (MMPs)} --- 27 MMPs with
      supporting components.
  \end{itemize}
  \begin{block}{Keep in mind}
    NeGP is the older umbrella \emph{plan}; e-Kranti later redefined it.
  \end{block}
\end{frame}

\begin{frame}{NeGD --- National e-Governance Division}
  \begin{itemize}
    \item \key{NeGD} = National e-Governance Division.
    \item Created in \key{2009} by the Ministry of Electronics and
      Information Technology (MeitY).
    \item An independent business division under the Digital India
      Corporation.
    \item Role: \key{programme management} and implementation support for
      e-Governance projects, with technical and advisory support.
  \end{itemize}
  \begin{block}{What it is}
    NeGD is an \emph{organisation} that supports e-Governance; it is not a
    plan.
  \end{block}
\end{frame}

\begin{frame}{NeGP vs e-Kranti vs NeGD}
  \begin{center}
    \begin{tabular}{p{0.13\textwidth}p{0.42\textwidth}p{0.11\textwidth}p{0.22\textwidth}}
      \toprule
      \textbf{Name} & \textbf{What it is} & \textbf{Year} & \textbf{One word} \\
      \midrule
      NeGP & The umbrella e-Governance plan (27 MMPs) & 2006 & The plan \\
      e-Kranti & Pillar 5 of Digital India; NeGP 2.0 & 2015 & The approach \\
      NeGD & Division under MeitY supporting projects & 2009 & The body \\
      \bottomrule
    \end{tabular}
  \end{center}
  \vspace{0.3cm}
  \begin{alertblock}{Trap alert}
    NeGD is a division, NeGP is a plan, e-Kranti is a pillar. Do not swap the
    three.
  \end{alertblock}
\end{frame}

\begin{frame}{e-NAM --- A Named Platform}
  \begin{itemize}
    \item \key{e-NAM} = National Agriculture Market (also written eNAM).
    \item A pan-India \key{electronic trading portal} for agricultural
      commodities.
    \item Launched on \key{14 April 2016}.
    \item Implemented by the \key{Small Farmers Agribusiness Consortium
      (SFAC)} under the Ministry of Agriculture and Farmers Welfare.
    \item Portal: \key{enam.gov.in}.
  \end{itemize}
\end{frame}

\begin{frame}{e-NAM and Agricultural Marketing}
  \begin{itemize}
    \item e-NAM \key{networks the existing APMC mandis} (Agricultural Produce
      Market Committee markets).
    \item It creates a \key{unified national market} for agricultural
      produce.
    \item Farmers get \key{online competitive bidding and transparent price
      discovery}; direct sale reduces intermediaries.
    \item \muted{AGMARKNET is a separate portal that gives agricultural
      marketing \emph{information} (market prices); e-NAM is the
      \emph{trading} platform.}
  \end{itemize}
  \begin{alertblock}{Trap alert}
    e-NAM is about agriculture marketing, not general e-commerce.
  \end{alertblock}
\end{frame}

\begin{frame}{The Citizen Portal}
  \begin{itemize}
    \item The \key{National Portal of India} is \key{india.gov.in}.
    \item It is a \key{single-window} gateway to information and services of
      the Government of India.
    \item \key{MyGov} is the citizen-engagement platform for participative
      governance.
  \end{itemize}
  \begin{block}{Single window}
    One entry point from which a citizen can reach many government services.
  \end{block}
\end{frame}

\begin{frame}{Named Platform --- Purpose Matching}
  \begin{center}
    \begin{tabular}{p{0.18\textwidth}p{0.66\textwidth}}
      \toprule
      \textbf{Platform} & \textbf{Purpose} \\
      \midrule
      e-NAM & Online trading of agricultural commodities \\
      GeM & Online government procurement of goods and services \\
      DigiLocker & Digital wallet for authentic documents \\
      UMANG & Single mobile app for many government services \\
      MyGov & Citizen engagement and participative governance \\
      e-Office & Paperless office and file management \\
      \bottomrule
    \end{tabular}
  \end{center}
  \vspace{0.2cm}
  \muted{Match the platform to its purpose, not to a similar-sounding name.}
\end{frame}

\begin{frame}{Electronic Records}
  \begin{itemize}
    \item The \key{Information Technology Act, 2000} gives legal recognition
      to \key{electronic records}.
    \item An electronic record can be stored, retrieved, and used in
      electronic form.
    \item A \key{digital signature} is used to authenticate an electronic
      record.
    \item \muted{Electronic records reduce paper, but they still need valid
      authentication.}
  \end{itemize}
  \begin{block}{Legal basis}
    It is the IT Act, 2000 that recognises electronic records and digital
    signatures.
  \end{block}
\end{frame}

\begin{frame}{Authentication}
  \begin{itemize}
    \item \key{Authentication} proves that a message or document is genuine
      and comes from the claimed sender.
    \item A \key{digital signature} provides \key{authenticity, integrity,
      and non-repudiation}.
    \item \key{Encryption} provides \key{confidentiality} (it keeps the
      content secret).
    \item \key{Aadhaar} is used as a digital identity for authentication of a
      person.
  \end{itemize}
  \begin{alertblock}{Trap alert}
    A digital signature is \emph{not} the same as encryption. Signature =
    authenticity; encryption = confidentiality (TRAP-15).
  \end{alertblock}
\end{frame}

\begin{frame}{Trap Alert: Confusions to Avoid}
  \begin{itemize}
    \item e-NAM = agricultural marketing; not general e-commerce.
    \item GeM = G2B procurement; not a G2C service.
    \item e-Office and PFMS = G2G; not G2C.
    \item NeGP (plan, 2006), e-Kranti (pillar/NeGP 2.0, 2015), NeGD
      (division, 2009) are three different things.
    \item A digital signature is not encryption.
  \end{itemize}
\end{frame}

\begin{frame}{Lesson 57: Recall}
  \begin{itemize}
    \item e-Governance uses ICT to deliver government services.
    \item G2x: G2C (citizen), G2G (government), G2B (business), G2E
      (employee).
    \item Digital India (1 July 2015): three vision areas, nine pillars;
      e-Kranti is Pillar 5.
    \item NeGP (2006, plan, 27 MMPs); e-Kranti (NeGP 2.0, pillar); NeGD
      (2009, MeitY division).
    \item e-NAM (14 April 2016) networks APMC mandis for agricultural
      marketing (SFAC).
    \item Citizen portal = india.gov.in; the IT Act, 2000 recognises
      electronic records; a digital signature provides authentication.
  \end{itemize}
\end{frame}
\end{document}
